\section{CƠ SỞ LÝ THUYẾT}

Để xây dựng một hệ thống phát hiện mã độc hiệu quả trên nền tảng Flower, việc thấu hiểu cơ chế toán học của FL và kiến trúc phần mềm của Flower là điều kiện tiên quyết.

\subsection{Toán học của Học Liên kết (Federated Learning)}
Mục tiêu của FL là tối ưu hóa một mô hình toàn cục (global model) với tham số $w$ mà không cần truy cập trực tiếp vào dữ liệu cục bộ $D_k$ của $K$ khách hàng (clients). Bài toán tối ưu hóa tổng quát có thể được biểu diễn như sau:
$$\min_{w} F(w) = \sum_{k=1}^{K} p_k F_k(w)$$
Trong đó:
\begin{itemize}
    \item $K$ là tổng số lượng client tham gia.
    \item $p_k$ là trọng số của client thứ $k$, thường được xác định bởi tỷ lệ dữ liệu sở hữu: $p_k = \frac{n_k}{n}$, với $n_k = |D_k|$ và $n = \sum n_k$.
    \item $F_k(w)$ là hàm mất mát cục bộ (local loss function) của client $k$, thường được định nghĩa là trung bình lỗi dự đoán trên tập dữ liệu $D_k$.
\end{itemize}

Thuật toán nền tảng được sử dụng rộng rãi nhất trong FL và cũng là chiến lược mặc định của Flower là Federated Averaging (FedAvg) \cite{mcmahan2017communication}. Quy trình hoạt động của FedAvg bao gồm 4 bước lặp lại trong mỗi vòng huấn luyện (communication round) $t$:

\begin{enumerate}
    \item \textbf{Selection \& Distribution:} Máy chủ trung tâm (Server) chọn một tập con các client $S_t$ (với tỷ lệ $C$) và gửi tham số mô hình toàn cục hiện tại $w_t$ cho họ.
    
    \item \textbf{Local Training:} Mỗi client $k \in S_t$ khởi tạo $w^k \leftarrow w_t$ và thực hiện cập nhật mô hình cục bộ bằng phương pháp Gradient Descent (SGD) trên dữ liệu riêng của mình qua $E$ epochs: \\
    $$w^k \leftarrow w^k - \eta \nabla F_k(w^k)$$ \\
    Sau quá trình này, client thu được bộ tham số mới $w_{t+1}^k$.
    
    \item \textbf{Upload:} Các client gửi tham số đã cập nhật $w_{t+1}^k$ (hoặc gradient $\Delta w_k = w_{t+1}^k - w_t$) trở lại máy chủ.
    
    \item \textbf{Aggregation:} Máy chủ tổng hợp các cập nhật để tạo ra mô hình toàn cục mới cho vòng tiếp theo: \\
    $$w_{t+1} = \sum_{k \in S_t} \frac{n_k}{n} w_{t+1}^k$$
\end{enumerate}
Sự ưu việt của Flower nằm ở chỗ nó trừu tượng hóa các bước toán học này thành các thành phần phần mềm có thể tùy chỉnh cao, cho phép các nhà nghiên cứu can thiệp vào bất kỳ bước nào, từ việc chọn client (sampling) đến công thức tổng hợp (aggregation strategy).

\subsection{Các Chiến lược Tổng hợp Bền vững (Robust Aggregation)}

Mặc dù FedAvg hoạt động hiệu quả trong môi trường lý tưởng, thuật toán này rất nhạy cảm với các cuộc tấn công đầu độc mô hình (Model Poisoning) hoặc dữ liệu nhiễu. Chỉ cần một client gửi về các tham số có giá trị cực lớn, mô hình toàn cục sẽ bị phá vỡ. Để khắc phục, hệ thống đề xuất sử dụng các cơ chế tổng hợp bền vững (Byzantine-robust aggregation) sau:

\subsubsection{Coordinate-wise Median (Trung vị theo tọa độ)}
Thay vì tính trung bình cộng, chiến lược này tính trung vị cho từng chiều của véc-tơ trọng số.
Với mỗi tham số thứ $j$ của mô hình, giá trị toàn cục mới được xác định bởi:
$$w_{j, t+1} = \text{median}(\{w_{j, t+1}^1, w_{j, t+1}^2, ..., w_{j, t+1}^K\})$$
Phương pháp này loại bỏ ảnh hưởng của các giá trị ngoại lai cực đoan, miễn là số lượng client độc hại chiếm ít hơn 50\% tổng số client.

\subsubsection{Trimmed Mean (Trung bình cắt cụt)}
Trimmed Mean là sự cân bằng giữa FedAvg và Median. Thuật toán này loại bỏ $m$ giá trị lớn nhất và $m$ giá trị nhỏ nhất ở mỗi chiều dữ liệu trước khi tính trung bình phần còn lại.
Công thức tổng quát cho tham số thứ $j$:
$$w_{j, t+1} = \frac{1}{|U_j|} \sum_{k \in U_j} w_{j, t+1}^k$$
Trong đó $U_j$ là tập hợp các chỉ số client còn lại sau khi đã loại bỏ $\beta$ tỷ lệ phần trăm các giá trị biên. Phương pháp này đặc biệt hiệu quả trong việc chống lại các tấn công làm lệch phân phối (distributional skew) trong khi vẫn giữ được tốc độ hội tụ tốt hơn Median.

\subsubsection{Krum và Multi-Krum}
Khác với hai phương pháp trên thực hiện trên từng tọa độ (coordinate-wise), Krum \cite{blanchard2017machine} lựa chọn một cập nhật mô hình từ một client thực tế mà "giống" với số đông nhất dựa trên khoảng cách Euclidean.
Với mỗi client $k$, ta tính điểm số $S(k)$ là tổng bình phương khoảng cách đến $K - f - 2$ client gần nhất ($f$ là số lượng client độc hại dự kiến):
$$S(k) = \sum_{i \rightarrow k} ||w^k - w^i||^2$$
Krum sẽ chọn update $w^{k^*}$ sao cho $S(k^*)$ là nhỏ nhất:
$$w_{t+1} = w_{t+1}^{k^*}$$
Krum đảm bảo tính hội tụ ngay cả khi có kẻ tấn công, tuy nhiên nhược điểm là tốc độ huấn luyện có thể chậm do chỉ sử dụng cập nhật từ một client duy nhất. Biến thể \textbf{Multi-Krum} khắc phục điều này bằng cách chọn ra $m$ client tốt nhất để tính trung bình thay vì chỉ một.

\subsection{Phân bố dữ liệu trong Học Liên kết (Data Partitioning)}

Một thách thức cốt lõi trong FL so với huấn luyện tập trung là tính chất thống kê của dữ liệu trên các client. Có hai kịch bản phân bố dữ liệu chính:

\subsubsection{IID (Independent and Identically Distributed)}
Trong trường hợp lý tưởng, dữ liệu trên mỗi client được coi là các mẫu độc lập và có cùng phân phối xác suất với dữ liệu tổng thể.
$$P_k(x, y) = P(x, y) \quad \forall k \in \{1, ..., K\}$$
Khi dữ liệu là IID, thuật toán FedAvg hội tụ nhanh và ổn định tương tự như huấn luyện tập trung.

\subsubsection{Non-IID (Dữ liệu không đồng nhất)}
Trong thực tế, đặc biệt là bài toán phát hiện mã độc, dữ liệu thường là Non-IID. Ví dụ: Một máy tính văn phòng chỉ nhiễm các loại macro virus, trong khi máy chủ lại nhiễm ransomware. Sự lệch pha này được chia thành:
\begin{itemize}
    \item \textbf{Label Distribution Skew (Lệch nhãn):} Client A chỉ có dữ liệu về mã độc loại A, Client B chỉ có loại B.
    \item \textbf{Feature Distribution Skew (Lệch đặc trưng):} Cùng một loại mã độc nhưng biến thể (signature) khác nhau trên các thiết bị khác nhau.
\end{itemize}
Dữ liệu Non-IID gây ra hiện tượng "Client Drift", nơi mô hình cục bộ của mỗi client hội tụ về các hướng khác nhau, khiến mô hình toàn cục sau khi tổng hợp bị giảm độ chính xác đáng kể.

\subsection{Kiến trúc Kỹ thuật của Flower Framework}

\subsubsection{Các Thành phần Cốt lõi}
Flower (hay flwr) được thiết kế với triết lý "client-agnostic" (không phụ thuộc client) và "framework-agnostic" (không phụ thuộc framework ML). Điều này có nghĩa là Flower có thể chạy trên các thiết bị biên yếu như Raspberry Pi, điện thoại di động, cho đến các cụm máy chủ lớn, và hỗ trợ PyTorch, TensorFlow, JAX hay thậm chí scikit-learn.

Hệ thống Flower hoạt động dựa trên mô hình Client-Server thông qua giao thức RPC (Remote Procedure Call).

\begin{itemize}
    \item \textbf{SuperLink (Bộ liên kết trung tâm):} Trong các phiên bản Flower hiện đại, SuperLink đóng vai trò trung gian quản lý kết nối. Nó duy trì trạng thái của hệ thống và điều phối các thông điệp giữa ServerApp và ClientApp.
    \item \textbf{ServerApp (Logic Máy chủ):} Đây là nơi chứa "trí tuệ" của quá trình FL. ServerApp thực thi Strategy, quyết định khi nào bắt đầu vòng mới, cấu hình tham số cho client, và xử lý kết quả trả về. ServerApp không trực tiếp thực hiện tính toán nặng (như lan truyền ngược) mà chỉ thực hiện các phép toán đại số tuyến tính trên các trọng số mô hình.
    \item \textbf{ClientApp (Logic Máy khách):} Chạy trên các nút biên. ClientApp đóng gói logic huấn luyện cục bộ. Nó nhận Config và Parameters từ Server, khởi tạo mô hình ML cục bộ (ví dụ: PyTorch Module), nạp dữ liệu từ ổ cứng, thực hiện huấn luyện (fit), và trả về kết quả.
    \item \textbf{Communication Stack (gRPC):} Flower sử dụng gRPC làm giao thức nền tảng để truyền tải các payload lớn (trọng số mô hình). gRPC sử dụng Protocol Buffers (Protobuf) để tuần tự hóa dữ liệu, giúp giảm kích thước gói tin và tăng tốc độ truyền tải so với REST/JSON truyền thống. Đặc biệt, gRPC hỗ trợ bảo mật TLS, cho phép xác thực máy chủ bảo mật cao - một yêu cầu quan trọng trong triển khai hệ thống an ninh mạng.
\end{itemize}

\subsubsection{Vòng đời của một Message trong Flower}
Hiểu rõ vòng đời message giúp chúng ta tùy chỉnh việc ghi log cho Dashboard Web:
\begin{enumerate}
    \item Server gọi \texttt{configure\_fit}: Strategy tạo ra các cấu hình (\texttt{FitIns}) cho các client được chọn.
    \item gRPC truyền \texttt{FitIns} (chứa weights toàn cục) đến Client.
    \item Client gọi \texttt{fit}: Thực hiện training cục bộ.
    \item Client trả về \texttt{FitRes}: Chứa weights mới và các metrics cục bộ (ví dụ: training loss, accuracy).
    \item Server gọi \texttt{aggregate\_fit}: Strategy nhận danh sách \texttt{FitRes}, tổng hợp trọng số mới, và quan trọng nhất cho dự án này, là tổng hợp metrics để ghi vào cơ sở dữ liệu giám sát.
\end{enumerate}

\subsection{Cơ chế Bảo mật kênh truyền thông (Secure Communication)}

Để đảm bảo tính toàn vẹn và bảo mật cho quá trình truyền tải trọng số mô hình (model weights) giữa Server và Client, hệ thống sử dụng cơ chế xác thực dựa trên Hạ tầng Khóa công khai (Public Key Infrastructure - PKI).

\subsubsection{SSL/TLS và Xác thực Phía Máy chủ (Server-side TLS)}
Hệ thống triển khai cơ chế \textbf{SSL/TLS phía máy chủ (Server-side TLS)}. Trong cấu hình này, máy chủ trình diện chứng chỉ số để máy khách xác thực danh tính trước khi thiết lập kết nối mã hóa, giúp ngăn chặn các máy chủ giả mạo.

\subsubsection{Quy trình cấp phát chứng chỉ (Certificates Workflow)}
Hệ thống sử dụng các thành phần sau để thiết lập niềm tin (Chain of Trust):
\begin{itemize}
    \item \textbf{Root CA (Certificate Authority):} Cơ quan cấp chứng chỉ gốc, do người quản trị hệ thống nắm giữ khóa bí mật. Đây là "nguồn gốc của niềm tin".
    \item \textbf{Server Certificate ($cert_S$):} Được ký bởi Root CA, dùng để chứng minh danh tính của ServerApp.
    \item \textbf{Client Certificate ($cert_C$):} Được ký bởi Root CA, cấp phát riêng cho từng thiết bị biên (Edge Device).
\end{itemize}

Khi kết nối gRPC được khởi tạo, một quá trình bắt tay (handshake) diễn ra. Client sử dụng chứng chỉ gốc Root CA để kiểm tra và xác thực $cert_S$ của Server. Nếu chứng chỉ hợp lệ, kênh truyền mã hóa (encrypted channel) mới được thiết lập. Mặc dù các chứng chỉ client đã được thiết lập sẵn, việc nâng cấp lên mTLS (xác thực hai chiều) được để lại như một hướng phát triển trong tương lai.

\subsection{Cơ sở lý thuyết về Quantum Machine Learning (QML)}

Quantum Machine Learning là sự giao thoa giữa Cơ học lượng tử và Học máy. Trong ngữ cảnh của đồ án này, chúng tôi tập trung vào việc sử dụng các Mạch lượng tử biến thiên (Variational Quantum Circuits - VQC) để thay thế hoặc bổ trợ cho các lớp mạng thần kinh cổ điển.

\subsubsection{Qubit và Nguyên lý chồng chập}
Khác với bit cổ điển chỉ nhận giá trị 0 hoặc 1, đơn vị cơ bản của thông tin lượng tử là Qubit, được biểu diễn dưới dạng véc-tơ trạng thái $|\psi\rangle$:
$$|\psi\rangle = \alpha|0\rangle + \beta|1\rangle$$
Trong đó $\alpha, \beta \in \mathbb{C}$ và $|\alpha|^2 + |\beta|^2 = 1$. Khả năng tồn tại đồng thời ở cả hai trạng thái (nguyên lý chồng chập) và hiện tượng vướng víu lượng tử (entanglement) cho phép QML xử lý không gian trạng thái lớn hơn theo cấp số nhân với số lượng tham số ít hơn.

\subsubsection{Quantum Convolutional Neural Network (QCNN)}
QCNN \cite{Cong_2019} là phiên bản lượng tử của CNN, được thiết kế để khai thác các đặc trưng cục bộ (local features) của dữ liệu. Cấu trúc của QCNN bao gồm ba thành phần chính:
\begin{enumerate}
    \item \textbf{Quantum Convolutional Layer:} Sử dụng các cổng lượng tử (quantum gates) tác động lên các qubit lân cận để trích xuất đặc trưng, tương tự như bộ lọc (filter) trong CNN.
    \item \textbf{Quantum Pooling Layer:} Đo lường một số qubit và sử dụng kết quả đó để điều khiển các qubit còn lại, giúp giảm chiều không gian Hilbert (tương tự như Pooling giảm kích thước ảnh).
    \item \textbf{Fully Connected Layer:} Lớp kết nối đầy đủ ở cuối để phân loại.
\end{enumerate}
Trong bài toán phát hiện mã độc, QCNN được kỳ vọng sẽ phát hiện được các mẫu (patterns) phức tạp trong chuỗi byte hoặc opcode mà CNN cổ điển có thể bỏ sót.

\subsubsection{Quantum Principal Component Analysis (QPCA)}
QPCA \cite{Lloyd_2014} là thuật toán giảm chiều dữ liệu dựa trên ma trận mật độ (density matrix).
Giả sử dữ liệu đầu vào là véc-tơ $u$, ta mã hóa nó thành trạng thái lượng tử $|\rho\rangle$. QPCA cho phép tìm ra các vector riêng (eigenvectors) và giá trị riêng (eigenvalues) của ma trận hiệp phương sai nhanh hơn so với thuật toán cổ điển theo hàm logarit của kích thước dữ liệu.
Trong hệ thống này, QPCA đóng vai trò tiền xử lý, giúp nén dữ liệu đặc trưng của mã độc trước khi đưa vào bộ phân loại, giúp giảm tải băng thông truyền tải trong mạng Federated Learning.

\subsubsection{Quantum Multilayer Perceptron (QMLP)}
QMLP \cite{shao2018quantummodelmultilayerperceptron} là kiến trúc mạng nơ-ron tổng quát sử dụng các cổng xoay tham số hóa (Parameterized Rotation Gates) như $R_x(\theta), R_y(\theta), R_z(\theta)$.
Quá trình huấn luyện QMLP thực chất là tìm bộ tham số $\theta$ tối ưu sao cho trạng thái đầu ra $|\psi_{out}\rangle$ gần với nhãn mục tiêu nhất. Ưu điểm của QMLP là khả năng biểu diễn (expressibility) cao hơn MLP cổ điển với cùng số lượng tham số, giúp mô hình nhẹ hơn và phù hợp để triển khai trên các thiết bị biên hạn chế tài nguyên.